% !TeX root = ../colloquium-2.tex

% Определение обратной функции.
% Критерий непрерывности монотонной функции.
% Теорема об обратной функции.

\phantomsection\label{sec_Definition04}
\begin{definition}[Определение обратной функции. 
                   Критерий непрерывности монотонной функции.
                   Теорема об обратной функции.]
    \subsection*{Определение обратной функции.}
    \begin{note}[purple]
        Пусть функция $f \; : \; E \rightarrow D$ осуществляет биекцию между $E$ и $D$.
        Если каждому $y \in D$ поставить в соответствие то $x \in E$, для которого $f(x) = y$, 
        то тем самым будет определена функция, отображающая множество $D$ во множество $E$. \\

        Эта функция будет называться \textbf{обратной} для функции $f$ и обозначается $f^{-1}$. \\
        Таким образом, $f^{-1} \; : \; D \rightarrow E$. 
    \end{note}

    \subsection*{Критерий непрерывности монотонной функции.}
    \begin{note}[purple]
        Монотонная на отрезке $[a, \, b]$ функция $f$ непрерывна на этом отрезке тогда и только тогда, 
        когда множеством её значений является отрезок с концами $f(a)$ и $f(b)$. 
    \end{note}

    \subsection*{Теорема об обратной функции.}
    \begin{note}[purple]
        Пусть функция $f$ непрерывна и строго монотонна (то есть возрастает или убывает) на 
        отрезке $[a, \, b]$. \\
        Тогда функция $f$ имеет обратную функцию $f^{-1}$, определённую на отрезке с концами 
        $f(a)$ и $f(b)$, причём $f^{-1}$ строго монотонна и непрерывна на отрезке с концами
        $f(a)$ и $f(b)$ и характер монотонности функций $f$ и $f^{-1}$ одинаковый. 
    \end{note}
\end{definition}